%% JOAS paper draft -- PRC Data Challenge 2026 (taxi-out time prediction)
%%
%% Template rules (full list in docs/paper/NOTES.md):
%%   - a SINGLE file: do NOT use \input{} or \include{}
%%   - do not rename the files: main.tex, figures/, reference.bib
%%   - title at most 12 words, Title Case
%%   - abstract a SINGLE paragraph, <=300 words, covering all four elements
%%   - write "Figure \ref{}", not "Fig."
%%   - use plain tabular; custom table styling breaks the HTML version
%%   - define an abbreviation only if it occurs more than 10 times in the text
%%
%% Every place marked [NUMBER] is filled in once the real data arrives.

\documentclass[
  manuscript=article,
  layout=preprint,
  year=2026,
  volume=x,
]{extra/joas}

\doi{xx.xxxxx/joas.xxxx.xxxx}

\received {x}
\revised  {x}
\accepted {x}
\published{x}

\addbibresource{reference.bib}

% Title: a 12 word limit. Candidate -- puts the contribution first, gives the airport count.
\title{Retrospective Taxi-Out Time Prediction at Eleven European Airports}

\author{Ahmet Abdullah Gultekin \orcid{0000-0000-0000-0000}}
\affiliation{Institution, City, Turkiye}
\email{correspondence@email.domain}

\keywords{taxi-out time; airport surface operations; departure queue; gradient boosting;
          post-operations analysis}

% An abbreviation only if it occurs more than 10 times. Review and prune once the writing is done.
\abbreviations{
    ATXOT: Additional Taxi-Out Time,
    RMSE: Root Mean Square Error
}

\begin{document}

\begin{abstract}
  % A SINGLE paragraph, <=300 words. Four elements are required:
  %   1) the purpose of the study and the problem investigated
  %   2) the basic design of the research approach
  %   3) the principal findings of the analysis
  %   4) a short summary of the interpretation and the conclusions
  % To be written last; a skeleton for now.
  Taxi-out time, the interval between off-block and take-off, is a primary driver of
  excess fuel burn on the airport surface, yet it is both hard to observe and hard to
  predict. [PURPOSE] This study predicts taxi-out time for departures at eleven major
  European airports using one full year of movement records. [DESIGN] We decompose the
  problem following EUROCONTROL's own additional taxi-out time indicator, reproducing its
  reference component and learning the residual with gradient boosting over congestion,
  routing and meteorological features. [FINDINGS] [NUMBER: RMSE, best configuration,
  family contributions] [INTERPRETATION] We further separate a retrospective model,
  appropriate for post-operations analysis, from a causal one usable in real time, and
  quantify the information value of retrospective observability. [NUMBER: the difference
  between the two models]
\end{abstract}


\section{Introduction}
% - why taxi-out matters: excess fuel and CO2, the PRC's indicator
% - why it is hard: a quantity that is hard to obtain (rationale.html), variability with
%   many causes
% - the question of this study: at 11 heterogeneous airports, WITHOUT trajectory data,
%   how well can it be predicted from movement records alone?
% - the contribution items (the three below), then the plan of the paper


\section{Related Work}
% To be developed from docs/literature.md. Three schools:
% - queueing theory: Idris et al. (2002) four factors, queue size the most important;
%   Simaiakis and Balakrishnan (2015) the decomposition
% - unimpeded time methods: Yin et al. (2017) criticise the percentile methods and
%   recommend regression -- our contribution sits in exactly that gap
% - European layout-based regression: Ravizza et al. (2013), ZRH is in our data set
% - ML applications: Herrema et al. (2018) a regression tree at LFPG, ~1.6 min;
%   Lee et al. (2016) departure fix; Wang et al. (2018) the SIFI/SCFI/AQLI/SRDI taxonomy
% Statement of the gap: all of them are single airport; nothing published on LTFM or LTAI.


\section{Data}
\subsection{Challenge Dataset}
% 11 airports, all of 2025, 4,167,797 movements; the ranking set is January + July 2026.
% For departures only the block time and the taxi time are blanked out -- describe the setup.
% [NUMBER] departures per airport, taxi-out distribution, missingness rates

\subsection{External Data}
% METAR (Iowa State IEM, public domain): [NUMBER] observations, 11 airports.
% OurAirports (public domain): coordinates (departure bearing), runway counts.
% Both are documented with their licence in docs/external_data.md.

\subsection{Data Quality}
% [NUMBER] does the identity MVT - BLOCK == TAXITIME hold
% [NUMBER] timestamp precision per airport (is any of it HH:MM)
% [NUMBER] the NM match rate


\section{Method}
\subsection{Reference Taxi-Out Time}
% A faithful reimplementation of the EUROCONTROL ATXOT: combination = (airport, stand,
% runway), reference = P10, validity = >=10 flights at or below P10, filters (>120 min,
% helicopters, de-icing after AOBT).
% Two deliberate deviations and the rationale for each:
%   1) rows without a valid reference cannot be dropped -> a hierarchical fallback chain
%   2) a fixed 2025 instead of a rolling 12 months
% [NUMBER] coverage at the official level, per airport

\subsection{Congestion Features}
% Idris's queue variable and the circularity problem; fixed-window proxies.
% Mapping onto the taxonomy of Wang et al. (SIFI/SCFI/AQLI/SRDI).
% The value of the arrivals not being blanked: the taxi-in median as a live congestion signal.

\subsection{Retrospective and Causal Anchoring}
% Two anchors of the same code. Why it is legitimate: the purpose of the competition is
% post-ops analysis.
% What the causal variant is good for: A-CDM, TSAT, DMAN.
% The two can be compared on the same validation set.

\subsection{Target Parameterisation}
% Raw taxi-out vs the residual over the ATXOT reference. Rationale and the [NUMBER] difference.

\subsection{Model and Validation}
% LightGBM, L2 objective (the optimal predictor under RMSE is the conditional mean), seed
% averaging.
% Validation: January and July held out from 2025 -- why random K-fold misleads.


\section{Results}
% MAIN TABLE: feature family x RMSE (the output of scripts/run_ablation.py).
% The per-airport RMSE table.
% The retrospective vs causal comparison.
% [NUMBER] everywhere.

\begin{table}[htbp!]
  \centering
  \small
  \caption{Feature family ablation. Negative delta means removing the family improves
           the model.}
  \label{tb:ablation}
  \begin{tabular}{lrrr}
  \toprule
  \textbf{Configuration} & \textbf{Features} & \textbf{RMSE (s)} & \textbf{$\Delta$} \\
  \midrule
  Full model & [NUMBER] & [NUMBER] & -- \\
  \bottomrule
  \end{tabular}
\end{table}


\section{Discussion}
% - which families carried and which did not (the negative results go here, openly)
% - why the airports differ: runway count (EDDM/EGLL 2, LTFM/EHAM 6), de-icing exposure
% - the January vs July difference and the de-icing regime
% - what we gain over the PRC's official indicator
% - limitations: no route or speed data, the reference does not roll, a single year


\section{Conclusion}


\section*{Acknowledgement}
% To EUROCONTROL PRC and the OpenSky Network for the data; to IEM and OurAirports for the
% open data.

\section*{Funding statement}
% If there is none, say so explicitly.

% The data statement is MANDATORY
\section*{Open data statement}
% Competition data set: PRC Data Challenge 2026, its DOI goes here once it is public.
% External data: IEM ASOS/METAR (public domain) and OurAirports (public domain) -- access URLs.
% Whether derived data can be shared: under the competition condition the 2026 set cannot be
% shared until it is public; that constraint will be stated explicitly here.

% The reproducibility statement is MANDATORY
\section*{Reproducibility statement}
% Source code: github.com/<user>/prc-taxiout-2026, GNU GPLv3.
% The code for the figures and tables: scripts/run_ablation.py, scripts/probe_data.py.
% Environment: pyproject.toml + uv.lock; no absolute paths, the data path comes from an
% environment variable.
% Tests: pytest tests/ ([NUMBER] tests).

\printbibliography

\end{document}
